\begin{PSbox}[unbreakable]{obj}{اهداف یادگیری}

پس از تکمیل این ماژول، دانشجویان قادر خواهند بود:

\begin{enumerate}
\def\labelenumi{\arabic{enumi}.}
\tightlist
\item
  توزیع‌های مشترک را به \lr{\mbox{N}} متغیر تصادفی تعمیم دهند
\item
  بردارهای تصادفی را تعریف کنند و با آن‌ها کار کنند
\item
  بردار میانگین و ماتریس‌های کوواریانس را محاسبه کنند
\item
  توزیع گاوسی چندمتغیره را به کار ببرند
\item
  تبدیل‌های خطی بردارهای تصادفی را انجام دهند
\end{enumerate}

\end{PSbox}

\PSsection{10.1}{مقدمه: از 2 به \lr{\mbox{N}} متغیر}

سیستم‌های واقعی اغلب شامل چندین کمیت نامعین به‌طور همزمان هستند:

\begin{itemize}
\tightlist
\item
  \textbf{آرایه حسگر:} \lr{\mbox{N}} حسگر که هر یک یک اندازه‌گیری نویزی تولید می‌کند
\item
  \textbf{سیستم \lr{\mbox{MIMO}}:} ضرایب کانال \lr{\mbox{M\ensuremath{\times}N}}
\item
  \textbf{برآورد پارامتر:} \lr{\mbox{p}} پارامتر نامعلوم
\item
  \textbf{نمونه‌های سیگنال:} \lr{\mbox{N}} نمونه نویز متوالی
\end{itemize}

به چارچوبی برای بردارهای تصادفی \textbf{\lr{\mbox{N}}-بعدی} نیاز داریم.

\PSsection{10.2}{توزیع‌های مشترک برای \lr{\mbox{N}} متغیر تصادفی}

\PSsubsection{\lr{\mbox{PDF}} مشترک}

برای \lr{\mbox{X\textsubscript{1}}}، \lr{\mbox{X\textsubscript{2}}}، \lr{…}، \lr{\mbox{X\textsubscript{n}}}:\PSdisplay{f_{X_1,...,X_n}(x_1,...,x_n) \geq 0, \quad \int\cdots\int f_{X_1,...,X_n} \, dx_1 \cdots dx_n = 1}

\PSsubsection{توزیع‌های حاشیه‌ای}

متغیرهای ناخواسته را انتگرال بگیرید (یا جمع کنید):\PSdisplay{f_{X_1}(x_1) = \int\cdots\int f_{X_1,...,X_n}(x_1,...,x_n) \, dx_2 \cdots dx_n}

\PSsubsection{توزیع‌های شرطی}

\PSdisplay{f_{X_1|X_2,...,X_n}(x_1|x_2,...,x_n) = \frac{f_{X_1,...,X_n}(x_1,...,x_n)}{f_{X_2,...,X_n}(x_2,...,x_n)}}

\PSsubsection{استقلال متقابل}

\lr{\mbox{X\textsubscript{1}}}، \lr{…}، \lr{\mbox{X\textsubscript{n}}} به‌طور متقابل مستقل هستند اگر:\PSdisplay{f_{X_1,...,X_n}(x_1,...,x_n) = \prod_{i=1}^n f_{X_i}(x_i)}

\PSsection{10.3}{بردارهای تصادفی}

\begin{PSbox}[unbreakable]{def}{تعریف}

یک \textbf{بردار تصادفی} مجموعه‌ای مرتب از متغیرهای تصادفی است:

\PSdisplay{\mathbf{X} = \begin{bmatrix} X_1 \\ X_2 \\ \vdots \\ X_n \end{bmatrix}}

\end{PSbox}

\PSsubsection{مثال‌های مهندسی}

{\def\LTcaptype{none} % do not increment counter
\begin{longtable}[c]{@{}ccc@{}}
\toprule\noalign{}
\rowcolor{PSnavy}\PSth{بردار تصادفی} & \PSth{مؤلفه‌ها} & \PSth{زمینه} \\
\midrule\noalign{}
\endhead
\bottomrule\noalign{}
\endlastfoot
خروجی آرایه حسگر & \lr{\mbox{X\textsubscript{1},…,X\textsubscript{n}}} = قرایت‌های حسگر & پردازش آرایه \\
کانال \lr{\mbox{MIMO}} & \lr{\mbox{h\textsubscript{1},…,h\textsubscript{n}}} = بهره‌های کانال & مخابرات بی‌سیم \\
بردار نویز & \lr{\mbox{N\textsubscript{1},…,N\textsubscript{n}}} = نمونه‌های نویز & پردازش سیگنال \\
بردار حالت & مکان، سرعت، شتاب & فیلتر کالمن \\
بردار ویژگی & اندازه‌گیری‌های استخراج‌شده & بازشناسی الگو \\
\end{longtable}
}

\PSsection{10.4}{بردار میانگین}

\begin{PSbox}[unbreakable]{def}{تعریف}

\PSdisplay{\boldsymbol{\mu} = E[\mathbf{X}] = \begin{bmatrix} E[X_1] \\ E[X_2] \\ \vdots \\ E[X_n] \end{bmatrix} = \begin{bmatrix} \mu_1 \\ \mu_2 \\ \vdots \\ \mu_n \end{bmatrix}}

\end{PSbox}

\PSsubsection{ویژگی‌ها}

\begin{itemize}
\tightlist
\item
  \lr{\mbox{E[AX + b] = AE[X] + b = A\ensuremath{\mu} + b}}
\item
  خطی‌بودن به بردارها و ماتریس‌ها تعمیم می‌یابد
\end{itemize}

\PSsection{10.5}{ماتریس کوواریانس}

\begin{PSbox}[unbreakable]{def}{تعریف}

\PSdisplay{\boldsymbol{\Sigma} = E[(\mathbf{X} - \boldsymbol{\mu})(\mathbf{X} - \boldsymbol{\mu})^T]}

درایه \lr{\mbox{(i,j)}}: \lr{\ensuremath{\Sigma}\textsubscript{ij} = Cov(X\textsubscript{i}, X\textsubscript{j})}

\PSdisplay{\boldsymbol{\Sigma} = \begin{bmatrix} \sigma_1^2 & \PSmtext{Cov}(X_1,X_2) & \cdots & \PSmtext{Cov}(X_1,X_n) \\ \PSmtext{Cov}(X_2,X_1) & \sigma_2^2 & \cdots & \PSmtext{Cov}(X_2,X_n) \\ \vdots & & \ddots & \vdots \\ \PSmtext{Cov}(X_n,X_1) & \cdots & \cdots & \sigma_n^2 \end{bmatrix}}

\end{PSbox}

\PSsubsection{ویژگی‌ها}

\begin{enumerate}
\def\labelenumi{\arabic{enumi}.}
\tightlist
\item
  \textbf{متقارن:} \lr{\ensuremath{\Sigma} = \ensuremath{\Sigma}\textsuperscript{T}}
\item
  \textbf{نیم‌معین مثبت:} \lr{\mbox{v\textsuperscript{T}\ensuremath{\Sigma}v \ensuremath{\ge} 0}} برای تمام \lr{\mbox{v}}
\item
  \textbf{قطر = واریانس‌ها:} \lr{\ensuremath{\Sigma}\textsubscript{ii} = Var(X\textsubscript{i})}
\item
  \textbf{غیرقطری = کوواریانس‌ها:} \lr{\ensuremath{\Sigma}\textsubscript{ij} = Cov(X\textsubscript{i}, X\textsubscript{j})}
\item
  \textbf{مؤلفه‌های مستقل \lr{\mbox{\ensuremath{\to} \ensuremath{\Sigma}}} قطری}
\end{enumerate}

\PSsubsection{فرمول جایگزین}

\PSdisplay{\boldsymbol{\Sigma} = E[\mathbf{X}\mathbf{X}^T] - \boldsymbol{\mu}\boldsymbol{\mu}^T}

\PSsection{10.6}{توزیع گاوسی چندمتغیره}

\begin{PSbox}[unbreakable]{def}{تعریف}

\lr{X \textasciitilde{} N(\ensuremath{\mu}, \ensuremath{\Sigma})} دارای \lr{\mbox{PDF}} زیر است:

\PSdisplay{f_\mathbf{X}(\mathbf{x}) = \frac{1}{(2\pi)^{n/2}|\boldsymbol{\Sigma}|^{1/2}} \exp\left(-\frac{1}{2}(\mathbf{x}-\boldsymbol{\mu})^T \boldsymbol{\Sigma}^{-1}(\mathbf{x}-\boldsymbol{\mu})\right)}

که در آن \lr{\mbox{\textbar{}\ensuremath{\Sigma}\textbar{}}} = دترمینان \lr{\mbox{\ensuremath{\Sigma}}}.

\end{PSbox}

\PSsubsection{ویژگی‌های کلیدی}

\begin{enumerate}
\def\labelenumi{\arabic{enumi}.}
\tightlist
\item
  \textbf{حاشیه‌ای‌ها گاوسی هستند:} هر زیرمجموعه‌ای از مؤلفه‌ها به‌طور مشترک گاوسی است
\item
  \textbf{شرطی‌ها گاوسی هستند:} \lr{X\textsubscript{1}\textbar{}X\textsubscript{2}=x\textsubscript{2}} گاوسی است
\item
  \textbf{تبدیل‌های خطی:} \lr{Y = AX + b \textasciitilde{} N(A\ensuremath{\mu}+b, A\ensuremath{\Sigma}A\textsuperscript{T})}
\item
  \textbf{ناهمبسته = مستقل} (ویژه گاوسی!)
\item
  \textbf{کاملاً مشخص‌شده} با \lr{\mbox{\ensuremath{\mu}}} و \lr{\mbox{\ensuremath{\Sigma}}} (فقط 2 پارامتر برای هر متغیر به‌علاوه همبستگی‌ها لازم است)
\end{enumerate}

\PSsubsection{چرا این‌قدر مهم است؟}

\begin{itemize}
\tightlist
\item
  قضیه حد مرکزی: جمع‌های بسیاری از متغیرهای تصادفی \lr{\mbox{\ensuremath{\to}}} گاوسی چندمتغیره
\item
  بردارهای نویز حرارتی گاوسی چندمتغیره هستند
\item
  از نظر تحلیلی خوش‌رفتار است (حاشیه‌ای‌ها و شرطی‌های فرم بسته)
\item
  فیلترهای بهینه (کالمن، وینر) گاوسی‌بودن را فرض می‌کنند
\end{itemize}

\PSsubsection{حالت دوبعدی (گاوسی دوبعدی)}

برای \lr{\mbox{n=2}} با \lr{\mbox{\ensuremath{\mu} = [0,0]\textsuperscript{T}}}، \lr{\ensuremath{\sigma}\textsubscript{1} = \ensuremath{\sigma}\textsubscript{2} = 1}:

\PSdisplay{f(x,y) = \frac{1}{2\pi\sqrt{1-\rho^2}} \exp\left(-\frac{x^2 - 2\rho xy + y^2}{2(1-\rho^2)}\right)}

خطوط تراز \textbf{بیضی} هستند (در حالت \lr{\mbox{\ensuremath{\rho} \ensuremath{\neq} 0}} کج شده‌اند).

\PSsection{10.7}{تبدیل‌های خطی بردارهای تصادفی}

\PSsubsection{تبدیل \lr{\mbox{Y = AX + b}}}

اگر \lr{\mbox{X}} یک بردار تصادفی با میانگین \lr{\mbox{\ensuremath{\mu}\textsubscript{X}}} و کوواریانس \lr{\mbox{\ensuremath{\Sigma}\textsubscript{X}}} باشد:

\PSdisplay{E[\mathbf{Y}] = A\boldsymbol{\mu}_X + \mathbf{b}}\PSdisplay{\boldsymbol{\Sigma}_Y = A\boldsymbol{\Sigma}_X A^T}

\PSsubsection{کاربردهای مهندسی}

\textbf{شکل‌دهی پرتو \lr{\mbox{(Beamforming)}}:} خروجی \lr{\mbox{y = w\textsuperscript{T}X}} (بردار وزن اعمال‌شده به آرایه حسگر)

\begin{itemize}
\tightlist
\item
  \lr{\mbox{E[y] = w\textsuperscript{T}\ensuremath{\mu}}}
\item
  \lr{\mbox{Var(y) = w\textsuperscript{T}\ensuremath{\Sigma}w}}
\end{itemize}

\textbf{سفیدسازی:} ماتریس \lr{\mbox{W}} را بیابید که در آن \lr{\mbox{W\ensuremath{\Sigma}W\textsuperscript{T} = I}} (داده را ناهمبسته می‌کند)

\begin{itemize}
\tightlist
\item
  \lr{\mbox{W = \ensuremath{\Sigma}\textsuperscript{-1/2}}}
\end{itemize}

\textbf{تحلیل مؤلفه‌های اصلی:} چرخش برای قطری‌کردن \lr{\mbox{\ensuremath{\Sigma}}}

\PSsection{10.8}{مثال‌های \lr{\mbox{MATLAB}}}

\begin{PSbox}[unbreakable]{ex}{مثال 1: تولید گاوسی چندمتغیره}

\tcbset{PScodemode/.style={unbreakable}}

\begin{Shaded}
\begin{Highlighting}[]
\CommentTok{\%\% Generate and Visualize 2D Gaussian Random Vector}
\VariableTok{mu} \OperatorTok{=}\NormalTok{ [}\FloatTok{1}\OperatorTok{;} \FloatTok{2}\NormalTok{]}\OperatorTok{;}
\VariableTok{Sigma} \OperatorTok{=}\NormalTok{ [}\FloatTok{4}\OperatorTok{,} \FloatTok{2}\OperatorTok{;} \FloatTok{2}\OperatorTok{,} \FloatTok{3}\NormalTok{]}\OperatorTok{;}  \CommentTok{\% Covariance matrix}

\VariableTok{N} \OperatorTok{=} \FloatTok{5000}\OperatorTok{;}
\VariableTok{X} \OperatorTok{=} \VariableTok{mvnrnd}\NormalTok{(}\VariableTok{mu}\OperatorTok{\textquotesingle{},} \VariableTok{Sigma}\OperatorTok{,} \VariableTok{N}\NormalTok{)}\OperatorTok{;}

\VariableTok{figure}\OperatorTok{;}
\VariableTok{scatter}\NormalTok{(}\VariableTok{X}\NormalTok{(}\OperatorTok{:,}\FloatTok{1}\NormalTok{)}\OperatorTok{,} \VariableTok{X}\NormalTok{(}\OperatorTok{:,}\FloatTok{2}\NormalTok{)}\OperatorTok{,} \FloatTok{5}\OperatorTok{,} \SpecialStringTok{\textquotesingle{}b\textquotesingle{}}\OperatorTok{,} \SpecialStringTok{\textquotesingle{}.\textquotesingle{}}\NormalTok{)}\OperatorTok{;}
\VariableTok{hold} \VariableTok{on}\OperatorTok{;}
\VariableTok{plot}\NormalTok{(}\VariableTok{mu}\NormalTok{(}\FloatTok{1}\NormalTok{)}\OperatorTok{,} \VariableTok{mu}\NormalTok{(}\FloatTok{2}\NormalTok{)}\OperatorTok{,} \SpecialStringTok{\textquotesingle{}r+\textquotesingle{}}\OperatorTok{,} \SpecialStringTok{\textquotesingle{}MarkerSize\textquotesingle{}}\OperatorTok{,} \FloatTok{15}\OperatorTok{,} \SpecialStringTok{\textquotesingle{}LineWidth\textquotesingle{}}\OperatorTok{,} \FloatTok{3}\NormalTok{)}\OperatorTok{;}
\VariableTok{xlabel}\NormalTok{(}\SpecialStringTok{\textquotesingle{}X\_1\textquotesingle{}}\NormalTok{)}\OperatorTok{;} \VariableTok{ylabel}\NormalTok{(}\SpecialStringTok{\textquotesingle{}X\_2\textquotesingle{}}\NormalTok{)}\OperatorTok{;}
\VariableTok{title}\NormalTok{(}\SpecialStringTok{\textquotesingle{}Bivariate Gaussian Samples\textquotesingle{}}\NormalTok{)}\OperatorTok{;}
\VariableTok{axis} \VariableTok{equal}\OperatorTok{;} \VariableTok{grid} \VariableTok{on}\OperatorTok{;}

\CommentTok{\% Verify statistics}
\VariableTok{fprintf}\NormalTok{(}\SpecialStringTok{\textquotesingle{}Mean: [\%.2f, \%.2f] (theory: [\%.1f, \%.1f])\textbackslash{}n\textquotesingle{}}\OperatorTok{,} \VariableTok{mean}\NormalTok{(}\VariableTok{X}\NormalTok{)}\OperatorTok{,} \VariableTok{mu}\OperatorTok{\textquotesingle{}}\NormalTok{)}\OperatorTok{;}
\VariableTok{fprintf}\NormalTok{(}\SpecialStringTok{\textquotesingle{}Cov matrix:\textbackslash{}n\textquotesingle{}}\NormalTok{)}\OperatorTok{;} \VariableTok{disp}\NormalTok{(}\VariableTok{cov}\NormalTok{(}\VariableTok{X}\NormalTok{))}\OperatorTok{;}
\end{Highlighting}
\end{Shaded}

\end{PSbox}

\begin{PSbox}[unbreakable]{ex}{مثال 2: تبدیل خطی}

\tcbset{PScodemode/.style={unbreakable}}

\begin{Shaded}
\begin{Highlighting}[]
\CommentTok{\%\% Transform: Y = AX + b}
\VariableTok{mu\_x} \OperatorTok{=}\NormalTok{ [}\FloatTok{0}\OperatorTok{;} \FloatTok{0}\NormalTok{]}\OperatorTok{;} \VariableTok{Sigma\_x} \OperatorTok{=} \VariableTok{eye}\NormalTok{(}\FloatTok{2}\NormalTok{)}\OperatorTok{;}   \CommentTok{\% Standard normal}
\VariableTok{A} \OperatorTok{=}\NormalTok{ [}\FloatTok{2}\OperatorTok{,} \FloatTok{1}\OperatorTok{;} \OperatorTok{{-}}\FloatTok{1}\OperatorTok{,} \FloatTok{3}\NormalTok{]}\OperatorTok{;}  \VariableTok{b} \OperatorTok{=}\NormalTok{ [}\FloatTok{1}\OperatorTok{;} \OperatorTok{{-}}\FloatTok{2}\NormalTok{]}\OperatorTok{;}

\VariableTok{N} \OperatorTok{=} \FloatTok{10000}\OperatorTok{;}
\VariableTok{X} \OperatorTok{=} \VariableTok{mvnrnd}\NormalTok{(}\VariableTok{mu\_x}\OperatorTok{\textquotesingle{},} \VariableTok{Sigma\_x}\OperatorTok{,} \VariableTok{N}\NormalTok{)}\OperatorTok{\textquotesingle{};}
\VariableTok{Y} \OperatorTok{=} \VariableTok{A}\OperatorTok{*}\VariableTok{X} \OperatorTok{+} \VariableTok{b}\OperatorTok{;}

\CommentTok{\% Theoretical}
\VariableTok{mu\_y\_theory} \OperatorTok{=} \VariableTok{A}\OperatorTok{*}\VariableTok{mu\_x} \OperatorTok{+} \VariableTok{b}\OperatorTok{;}
\VariableTok{Sigma\_y\_theory} \OperatorTok{=} \VariableTok{A}\OperatorTok{*}\VariableTok{Sigma\_x}\OperatorTok{*}\VariableTok{A}\OperatorTok{\textquotesingle{};}

\VariableTok{fprintf}\NormalTok{(}\SpecialStringTok{\textquotesingle{}E[Y] theory: [\%.1f; \%.1f], sim: [\%.2f; \%.2f]\textbackslash{}n\textquotesingle{}}\OperatorTok{,} \OperatorTok{...}
    \VariableTok{mu\_y\_theory}\OperatorTok{,} \VariableTok{mean}\NormalTok{(}\VariableTok{Y}\OperatorTok{,}\FloatTok{2}\NormalTok{))}\OperatorTok{;}
\VariableTok{fprintf}\NormalTok{(}\SpecialStringTok{\textquotesingle{}Σ\_Y theory:\textbackslash{}n\textquotesingle{}}\NormalTok{)}\OperatorTok{;} \VariableTok{disp}\NormalTok{(}\VariableTok{Sigma\_y\_theory}\NormalTok{)}\OperatorTok{;}
\VariableTok{fprintf}\NormalTok{(}\SpecialStringTok{\textquotesingle{}Σ\_Y simulated:\textbackslash{}n\textquotesingle{}}\NormalTok{)}\OperatorTok{;} \VariableTok{disp}\NormalTok{(}\VariableTok{cov}\NormalTok{(}\VariableTok{Y}\OperatorTok{\textquotesingle{}}\NormalTok{))}\OperatorTok{;}
\end{Highlighting}
\end{Shaded}

\end{PSbox}

\begin{PSbox}[breakable]{ex}{مثال 3: آرایه حسگر با نویز همبسته}

\tcbset{PScodemode/.style={unbreakable}}

\begin{Shaded}
\begin{Highlighting}[]
\CommentTok{\%\% 4{-}Element Sensor Array with Spatially Correlated Noise}
\VariableTok{n\_sensors} \OperatorTok{=} \FloatTok{4}\OperatorTok{;}
\VariableTok{sigma2} \OperatorTok{=} \FloatTok{1}\OperatorTok{;}      \CommentTok{\% Noise power per sensor}
\VariableTok{rho} \OperatorTok{=} \FloatTok{0.6}\OperatorTok{;}       \CommentTok{\% Adjacent sensor correlation}

\CommentTok{\% Build covariance matrix (exponential decay model)}
\VariableTok{Sigma} \OperatorTok{=} \VariableTok{zeros}\NormalTok{(}\VariableTok{n\_sensors}\NormalTok{)}\OperatorTok{;}
\KeywordTok{for} \VariableTok{i} \OperatorTok{=} \FloatTok{1}\OperatorTok{:}\VariableTok{n\_sensors}
    \KeywordTok{for} \VariableTok{j} \OperatorTok{=} \FloatTok{1}\OperatorTok{:}\VariableTok{n\_sensors}
        \VariableTok{Sigma}\NormalTok{(}\VariableTok{i}\OperatorTok{,}\VariableTok{j}\NormalTok{) }\OperatorTok{=} \VariableTok{sigma2} \OperatorTok{*} \VariableTok{rho}\OperatorTok{\^{}}\VariableTok{abs}\NormalTok{(}\VariableTok{i}\OperatorTok{{-}}\VariableTok{j}\NormalTok{)}\OperatorTok{;}
    \KeywordTok{end}
\KeywordTok{end}

\VariableTok{fprintf}\NormalTok{(}\SpecialStringTok{\textquotesingle{}Noise covariance matrix:\textbackslash{}n\textquotesingle{}}\NormalTok{)}\OperatorTok{;} \VariableTok{disp}\NormalTok{(}\VariableTok{Sigma}\NormalTok{)}\OperatorTok{;}

\CommentTok{\% Generate noise samples}
\VariableTok{N\_samp} \OperatorTok{=} \FloatTok{10000}\OperatorTok{;}
\VariableTok{L} \OperatorTok{=} \VariableTok{chol}\NormalTok{(}\VariableTok{Sigma}\OperatorTok{,} \SpecialStringTok{\textquotesingle{}lower\textquotesingle{}}\NormalTok{)}\OperatorTok{;}
\VariableTok{noise} \OperatorTok{=} \VariableTok{L} \OperatorTok{*} \VariableTok{randn}\NormalTok{(}\VariableTok{n\_sensors}\OperatorTok{,} \VariableTok{N\_samp}\NormalTok{)}\OperatorTok{;}

\CommentTok{\% Signal present in direction: s = [1; 1; 1; 1] (broadside)}
\VariableTok{signal} \OperatorTok{=} \FloatTok{3}\OperatorTok{;}
\VariableTok{s} \OperatorTok{=} \VariableTok{ones}\NormalTok{(}\VariableTok{n\_sensors}\OperatorTok{,} \FloatTok{1}\NormalTok{) }\OperatorTok{/} \VariableTok{sqrt}\NormalTok{(}\VariableTok{n\_sensors}\NormalTok{)}\OperatorTok{;}
\VariableTok{received} \OperatorTok{=} \VariableTok{signal} \OperatorTok{*} \VariableTok{s} \OperatorTok{+} \VariableTok{noise}\OperatorTok{;}  \CommentTok{\% N\_sensors x N\_samp}

\CommentTok{\% Beamformer output: y = w\textquotesingle{}*received (w = s for matched filter)}
\VariableTok{w} \OperatorTok{=} \VariableTok{s}\OperatorTok{;}
\VariableTok{y} \OperatorTok{=} \VariableTok{w}\OperatorTok{\textquotesingle{}} \OperatorTok{*} \VariableTok{received}\OperatorTok{;}
\VariableTok{SNR\_out} \OperatorTok{=} \VariableTok{signal}\OperatorTok{\^{}}\FloatTok{2} \OperatorTok{/}\NormalTok{ (}\VariableTok{w}\OperatorTok{\textquotesingle{}*}\VariableTok{Sigma}\OperatorTok{*}\VariableTok{w}\NormalTok{)}\OperatorTok{;}
\VariableTok{fprintf}\NormalTok{(}\SpecialStringTok{\textquotesingle{}Output SNR = \%.2f (\%.1f dB)\textbackslash{}n\textquotesingle{}}\OperatorTok{,} \VariableTok{SNR\_out}\OperatorTok{,} \FloatTok{10}\OperatorTok{*}\VariableTok{log10}\NormalTok{(}\VariableTok{SNR\_out}\NormalTok{))}\OperatorTok{;}
\end{Highlighting}
\end{Shaded}

\end{PSbox}

\PSsection{10.9}{مسایل تمرینی}

\begin{PSbox}[unbreakable]{prob}{مسیله 1}

بردار تصادفی \lr{X = [X\textsubscript{1}, X\textsubscript{2}, X\textsubscript{3}]\textsuperscript{T}} با \lr{\mbox{\ensuremath{\mu} = [1, 0, -1]\textsuperscript{T}}} و ماتریس کوواریانس \lr{\mbox{\ensuremath{\Sigma} = [4 1 0; 1 2 -1; 0 -1 3]}}.

\begin{enumerate}
\def\labelenumi{(\PSalph{enumi})}
\tightlist
\item
  \lr{Var(X\textsubscript{1} + X\textsubscript{2} + X\textsubscript{3})} را بیابید. (ب) \lr{Cov(X\textsubscript{1}, X\textsubscript{2}+X\textsubscript{3})} را بیابید. (ج) آیا \lr{\mbox{X\textsubscript{1}}} و \lr{\mbox{X\textsubscript{3}}} ناهمبسته هستند؟
\end{enumerate}

\PSsolution{} 

\begin{enumerate}
\def\labelenumi{(\PSalph{enumi})}
\tightlist
\item
  \lr{Var(X\textsubscript{1}+X\textsubscript{2}+X\textsubscript{3}) = 1\textsuperscript{T}\ensuremath{\Sigma}1 = 4+2+3 + 2(1) + 2(0) + 2(-1) = 9 + 0 = 9}\PSbr{}توجه: = \lr{\mbox{\ensuremath{\Sigma}\textsubscript{ij}}} تمام درایه‌ها = \lr{\mbox{4+1+0+1+2+(-1)+0+(-1)+3 = 9}}
\item
  \lr{Cov(X\textsubscript{1}, X\textsubscript{2}+X\textsubscript{3}) = Cov(X\textsubscript{1},X\textsubscript{2}) + Cov(X\textsubscript{1},X\textsubscript{3}) = 1 + 0 = 1}
\item
  \lr{\mbox{Cov(X\textsubscript{1},X\textsubscript{3}) = 0}}، بنابراین بله، \lr{\mbox{X\textsubscript{1}}} و \lr{\mbox{X\textsubscript{3}}} ناهمبسته هستند.
\end{enumerate}

\end{PSbox}

\begin{PSbox}[unbreakable]{prob}{مسیله 2}

\lr{\mbox{X \textasciitilde{} N(0, \ensuremath{\Sigma})}} با \lr{\ensuremath{\Sigma} = [1 \ensuremath{\rho}; \ensuremath{\rho} 1]}. توزیع شرطی \lr{X\textsubscript{1}\textbar{}X\textsubscript{2}=x\textsubscript{2}} را بیابید.

\PSsolution{} برای گاوسی دوبعدی:\PSbr{}\lr{X\textsubscript{1}\textbar{}X\textsubscript{2}=x\textsubscript{2} \textasciitilde{} N(\ensuremath{\mu}\textsubscript{1} + \ensuremath{\rho}(\ensuremath{\sigma}\textsubscript{1}/\ensuremath{\sigma}\textsubscript{2})(x\textsubscript{2}-\ensuremath{\mu}\textsubscript{2}), \ensuremath{\sigma}\textsubscript{1}\textsuperscript{2}(1-\ensuremath{\rho}\textsuperscript{2}))}\PSbr{}\lr{= N(\ensuremath{\rho}x\textsubscript{2}, 1-\ensuremath{\rho}\textsuperscript{2})}

میانگین شرطی تابعی خطی از \lr{\mbox{x\textsubscript{2}}} است و واریانس شرطی با ضریب \lr{\mbox{(1-\ensuremath{\rho}\textsuperscript{2})}} کاهش می‌یابد.

\end{PSbox}

\begin{PSbox}[unbreakable]{prob}{مسیله 3}

کد \lr{\mbox{MATLAB}} را برای تولید یک بردار گاوسی سه‌بعدی با \lr{\mbox{\ensuremath{\mu}}} و \lr{\mbox{\ensuremath{\Sigma}}} معلوم، محاسبه کوواریانس نمونه و بررسی آن بنویسید.

\PSsolution{} 

\tcbset{PScodemode/.style={unbreakable}}

\begin{Shaded}
\begin{Highlighting}[]
\VariableTok{mu} \OperatorTok{=}\NormalTok{ [}\FloatTok{1}\OperatorTok{;} \FloatTok{0}\OperatorTok{;} \OperatorTok{{-}}\FloatTok{1}\NormalTok{]}\OperatorTok{;}
\VariableTok{Sigma} \OperatorTok{=}\NormalTok{ [}\FloatTok{4} \FloatTok{1} \FloatTok{0}\OperatorTok{;} \FloatTok{1} \FloatTok{2} \OperatorTok{{-}}\FloatTok{1}\OperatorTok{;} \FloatTok{0} \OperatorTok{{-}}\FloatTok{1} \FloatTok{3}\NormalTok{]}\OperatorTok{;}
\VariableTok{N} \OperatorTok{=} \FloatTok{50000}\OperatorTok{;}
\VariableTok{X} \OperatorTok{=} \VariableTok{mvnrnd}\NormalTok{(}\VariableTok{mu}\OperatorTok{\textquotesingle{},} \VariableTok{Sigma}\OperatorTok{,} \VariableTok{N}\NormalTok{)}\OperatorTok{;}
\VariableTok{fprintf}\NormalTok{(}\SpecialStringTok{\textquotesingle{}Sample mean: \textquotesingle{}}\NormalTok{)}\OperatorTok{;} \VariableTok{disp}\NormalTok{(}\VariableTok{mean}\NormalTok{(}\VariableTok{X}\NormalTok{)}\OperatorTok{\textquotesingle{}}\NormalTok{)}\OperatorTok{;}
\VariableTok{fprintf}\NormalTok{(}\SpecialStringTok{\textquotesingle{}Sample cov:\textbackslash{}n\textquotesingle{}}\NormalTok{)}\OperatorTok{;} \VariableTok{disp}\NormalTok{(}\VariableTok{cov}\NormalTok{(}\VariableTok{X}\NormalTok{))}\OperatorTok{;}
\end{Highlighting}
\end{Shaded}

\emph{ماژول بعدی: {ماژول 11 \lr{—} قضیه حد مرکزی}}

\end{PSbox}
